\documentclass[aspectratio=169]{beamer}
\usetheme{metropolis}

\usepackage{amsmath}
\usepackage{mathdots}
\usepackage{mathtools}
\usepackage{amssymb}
\usepackage{xcolor}
\usepackage{tipa}
\usefonttheme[onlymath]{serif}

\definecolor{sifivered}{HTML}{FC3F4E}
\definecolor{sifivered2}{HTML}{F86C66}
\definecolor{sifivegrey}{HTML}{32393E}
\definecolor{sifivegrey2}{HTML}{212629}
\definecolor{sifiveyellow}{HTML}{FFCC00}

\setbeamercolor{progress bar}{fg=sifivegrey}
\setbeamercolor{title separator}{fg=sifivegrey}
\setbeamercolor{progress bar in head/foot}{fg=sifivegrey}
\setbeamercolor{progress bar in section page}{fg=sifivegrey}

\title{Informing Physical Design in SCD}
\subtitle{\textbf{Predicting parameter cost}}
\author{Robert D'Angelo}
%\institute{SiFive}
\date{\today}

%\usebackgroundtemplate%
%{%
    %\includegraphics[width=\paperwidth,height=\paperheight]{./figures/sifive_presentation_background_alpha.pdf}%
%}

\defbeamertemplate{frame numbering}{altcounter}{
\vspace{-2.5mm}
$|$\insertframenumber
\hspace{7mm}
}
\setbeamertemplate{frame numbering}[altcounter]
\setbeamercolor{frame numbering}{fg=white}

\setbeamercolor{section in head/foot}{fg=white, bg=sifivegrey}

\makeatletter
\setbeamertemplate{footline}
{
  \leavevmode%
  \hbox{%
  \begin{beamercolorbox}[wd=.333333\paperwidth,ht=2.25ex,dp=1ex,center]{section in head/foot}%
    \usebeamerfont{author in head/foot}\insertshortauthor ~-- \beamer@ifempty{\insertshortinstitute}{}{SiFive}
  \end{beamercolorbox}%
  \begin{beamercolorbox}[wd=.333333\paperwidth,ht=2.25ex,dp=1ex,center]{section in head/foot}%
    \usebeamerfont{title in head/foot}\insertshorttitle
  \end{beamercolorbox}%
  \begin{beamercolorbox}[wd=.333333\paperwidth,ht=2.25ex,dp=1ex,right]{section in head/foot}%
    \usebeamerfont{date in head/foot}\insertshortdate{}\hspace*{2em}
    \insertframenumber{} / \inserttotalframenumber\hspace*{2ex} 
  \end{beamercolorbox}}%
  \vskip0pt%
}
\makeatother

\newcommand\myfootnote[1]{%
  \tikz[remember picture,overlay]
  \draw (current page.south west) +(1in + \oddsidemargin,2em)
  node[anchor=south west,inner sep=0pt]{\parbox{\textwidth}{%
      \rlap{\rule{10em}{0.4pt}}\raggedright\tiny#1}};
}


% for schematics graphics
\newcommand{\trimmedgraphic}[1]{%
    %\centering
    \centerline{
        \makebox[\textwidth]{
                    \includegraphics[trim = 5cm 3cm 3cm 8.5cm,clip,width=0.9\linewidth]{#1}
                }
    }
}

\newcommand{\trimmedgraphicconv}[2]{
    \begin{figure}
        \centering
        \includegraphics[trim = 2cm 0cm 2cm 0cm,clip,width=0.33\linewidth]{#1}
        \caption*{#2}
    \end{figure}
}

% place graphic in the center with max size on the slide
\newcommand{\centeredgraphic}[1]{
    \centerline{
        \makebox[\textwidth]{
                    \includegraphics[width=\linewidth,height=0.8\textheight, keepaspectratio]{#1}
                }
    }
}

% place graphic in the center with max size on the slide
\newcommand{\centeredgraphicsmall}[1]{
    \centerline{
        \makebox[\textwidth]{
                    \includegraphics[height=0.5\textheight, keepaspectratio]{#1}
                }
    }
}


\newcommand{\centeredgraphictwo}[2]{
    \centerline{
        \makebox[\textwidth]{
                    \includegraphics[width=0.5\linewidth,height=0.8\textheight, keepaspectratio]{#1}
                    \includegraphics[width=0.5\linewidth,height=0.8\textheight, keepaspectratio]{#2}
                }
    }
}
\newcommand{\centeredgraphictwovert}[2]{
    \centerline{
        \makebox[\textwidth]{
                \begin{tabular}{c}
                    \includegraphics[height=0.4\textheight, keepaspectratio]{#1} \\
                    \includegraphics[height=0.4\textheight, keepaspectratio]{#2}
                \end{tabular}
                }
    }
}



% place graphic in the center with max size on the slide with transparency
\newcommand{\centeredgraphictransparent}[1]{
    \centerline{
        \makebox[\textwidth]{
                    \includegraphicstransparent[width=\linewidth,height=0.8\textheight, keepaspectratio]{#1}{15}
                }
    }
}

\newcommand{\centeredgraphicwidthtransparent}[2]{
    \centerline{
        \makebox[\textwidth]{
                    \includegraphicstransparent[width=#2\linewidth,height=0.8\textheight, keepaspectratio]{#1}{15}
                }
    }
}

% place graphic in the center with given width
\newcommand{\centeredgraphicwidth}[2]{
    \centerline{
        \makebox[\textwidth]{
                    \includegraphics[width=#2\linewidth,height=0.8\textheight, keepaspectratio]{#1}
                }
    }
}

% place graphic in the center with height and width specified, but aspect ratio maintained
\newcommand{\centeredgraphicaspect}[3]{
    \centerline{
        \makebox[\textwidth]{
                    \includegraphics[width=#2\linewidth,height=#3\textheight, keepaspectratio]{#1}
                }
    }
}

\newcommand{\centeredgraphicaspectfootnote}[4]{
    \centerline{
        \makebox[\textwidth]{
                    \includegraphics[width=#2\linewidth,height=#3\textheight, keepaspectratio]{#1}
                }
    }
    \myfootnote{#4}
}

% place a frame that is just an image and a title
\newcommand{\imageonlyframetwo}[4]{% image, image, title, width
\begin{frame}
    \frametitle{#3}
    \pause
    \begin{overprint}
    \onslide<1>\centeredgraphicwidth{#1}{#4}
    \onslide<2>\centeredgraphicwidthtransparent{#1}{#4}
    \end{overprint}

    \begin{overprint}
    \onslide<1>\centeredgraphicwidthtransparent{#2}{#4}
    \onslide<2>\centeredgraphicwidth{#2}{#4}
    \end{overprint}
\end{frame}
}

% place a frame that is just an image and a title
\newcommand{\imageonlyframe}[3]{% image, title, caption
\begin{frame}
    \frametitle{#2}
    #3
    \visible<1>{\centeredgraphic{#1}}
\end{frame}
}

\newcommand{\imageonlyframenonum}[3]{% image, title, caption
\begin{frame}[noframenumbering]
    \frametitle{#2}
    #3
    \visible<1>{\centeredgraphic{#1}}
\end{frame}
}



\newcommand{\imageonlyframetest}[4]{% image, title, caption
\begin{frame}
    \frametitle{#3 \includegraphics[width=0.05\linewidth, keepaspectratio]{#2}}
    #4
    \visible<1>{\centeredgraphic{#1}}
\end{frame}
}


% place a frame that is just an image and a title
\newcommand{\imageonlyframenopause}[3]{% image, title, caption
\begin{frame}
    \frametitle{#2}
    #3
    \centeredgraphic{#1}
\end{frame}
}


%\newcommand\blfootnote[1]{%
  %\begingroup
  %\renewcommand\thefootnote{}\footnote{#1}%
  %\addtocounter{footnote}{-1}%
  %\endgroup
%}

%\newcommand{\myfootnote}[1]{
    %\setbeamerfont{footnote}{size=\tiny}
    %\blfootnote{#1}
%}

% place a frame that is just an image and a title and a footnote
\newcommand{\imageonlyframefootnote}[4]{% image, title, footnote, caption
\begin{frame}
    \frametitle{#2}
    #4
    \vspace{-0.5cm}
    \centeredgraphic{#1}
    \myfootnote{#3}
\end{frame}
}

\newcommand{\imageonlyframefootnotenoskip}[4]{% image, title, footnote, caption
\begin{frame}
    \frametitle{#2}
    #4
    \centeredgraphic{#1}
    \myfootnote{#3}
\end{frame}
}

\newcommand{\imageonlyframefootnotenoskipnonum}[4]{% image, title, footnote, caption
\begin{frame}[noframenumbering]
    \frametitle{#2}
    #4
    \centeredgraphic{#1}
    \myfootnote{#3}
\end{frame}
}



\newcommand{\imageonlyframefootnoteaspect}[6]{% image, title, footnote, caption, width, height
\begin{frame}
    \frametitle{#2}
    #4
    \centeredgraphicaspect{#1}{#5}{#6}
    \myfootnote{#3}
\end{frame}
}






\begin{document}

%\begin{frame}
    %\titlepage{}
%\end{frame}

%\begin{frame}
    %\frametitle{Overview}
    %\tableofcontents{}

    %\section{Analog computation}
        %\subsection{Analog multipliers}
        %\subsection{Time-mode arithmetic circuits}
    %\subsection{Scalable analog circuits}
        %\subsection{Time-mode circuis}
        %\subsection{Bulk driven circuits}
    %\section{Analog to information converters}
        %\subsection{Compressed sensing ADCs}
        %\subsection{Machine learning ADCs}
    %\section{Biomedical signals}
        %\subsection{ExG: EEG, ECG, EOG, etc.}
%\end{frame}

\bgroup
\makeatletter
\setbeamertemplate{footline}
{
  \leavevmode%
  \hbox{%
  %
  %
  }%
  \vskip0pt%
}
\makeatother

\setbeamercolor{normal text}{fg=white}
\setbeamercolor{frametitle}{fg=white}
\usebeamercolor[fg]{normal text}
{
\usebackgroundtemplate{\includegraphics[width=\paperwidth]{./figures/sifive_presentation_title_page.pdf}}%

\begin{frame}
    \titlepage{}
    %\centerline{\includegraphics[scale=0.5]{figures/source_mixing_cichocki2005.pdf}}
\end{frame}
}
\egroup

\setcounter{framenumber}{0}

\setbeamercolor{normal text}{fg=black}
\setbeamercolor{frametitle}{fg=sifivered, bg=sifivegrey}
%\usebeamercolor[fg]{normal text}
%\begin{frame}
    %\frametitle{Charge pump topologies}
    %\begin{itemize}
    %\item Dickinson
        %\begin{itemize}
        %\item on chip caps
        %\item body effect issues
        %\end{itemize}
    %\end{itemize}
    %%\centerline{\includegraphics[scale=0.5]{figures/source_mixing_cichocki2005.pdf}}
%\end{frame}

\imageonlyframe{./figures/1200px-PhysicalDesign.png} % figure
               {What is physical design?}             % title
               {} % caption

\begin{frame}
    \frametitle{How can SCD help?}
    Real time feedback on “cost” of configuration choices.
    \begin{itemize}
        \item "Impact scores" per knob
        \begin{itemize}
            \item Area
            \item Power
            \item Timing
            \item Compute cost (PNR/timing optimization), e.g. clock gating
            \item I/O requirements
            \item Benchmarks 
        \end{itemize}
    \end{itemize}
\end{frame}

\begin{frame}
    \frametitle{Why show this information? Motivating example}
    Yield/Cost vs. Area
    %\vspace{2mm}
    \centeredgraphictwo{figures/point_defects.png}{figures/yield_area_cost.png}
    4x area reduction $\rightarrow$ 10x cost reduction
    %\begin{align}
       %V_{n} = \frac{C V_{cc}}{C + C_s} - V_{t}\left[V_{\scriptscriptstyle SB}^{n}\right]
    %\end{align}

    %\myfootnote{Pan, F., \& Samaddar, T. (2010). Charge pump circuit design. McGraw Hill Professional.}
\end{frame}

\imageonlyframe{./figures/gradient_descent.pdf} % figure
               {Real time knob cost predictions in your browser?!?} % title
               {} % caption

\begin{frame}
    \frametitle{Need a model}
    Challenges to making front end prediction
    \begin{itemize}
        \item Compute time (Too many parameters)
        \item Modeling (Complex designs)
        \item Licensing (Expensive software, prohibitive use)
    \end{itemize}
\end{frame}

\begin{frame}
    \frametitle{Need a model, Part 2}
    Single input, single output neural network may not be enough. Ideally, create a model that can map data from different correlated data sources to qualitative performance metrics.
    \begin{itemize}
        \item Compute time $\rightarrow$ Sparse sampling + Pattern recognition/ML
        \item Modeling $\rightarrow$ Statistical dependence over physical models
        \item Licensing $\rightarrow$ Remote execution and FOSS
    \end{itemize}
\end{frame}

\begin{frame}
    \frametitle{Need data}
    \begin{itemize}
        \item Hardware measurements
        \begin{itemize}
            \item Best, ground truth, fast (speed of RL)
            \item Expensive (\$10k-\$1M), Time consuming (months to years)
            \item Error prone (humans make mistakes), outlier effects (small sample size)
            \item Cutting edge prohibitively expensive ($>14nm$)
        \end{itemize}
    \end{itemize}
    \centeredgraphicsmall{figures/sifive_chip.jpg}
\end{frame}

\begin{frame}
    \frametitle{Need data}
    \begin{itemize}
        \item ASIC Simulations (SPICE, PNR annotated gate sim, .lib delay models)
        \begin{itemize}
            \item Accurate, time consuming (hours to days to weeks)
            \item Can do some statistics for free with stat params from PDK
            \item Expensive (\$40k per license per year), time (days to weeks, 10 hrs with prometheus) 
            \item Cutting edge ($<14nm$ simulations)
        \end{itemize}
    \end{itemize}
        \centeredgraphicsmall{figures/td_wave.png}
\end{frame}
\begin{frame}
    \frametitle{Need data}
    \begin{itemize}
        \item FPGA 
        \begin{itemize}
            \item Apples vs. oranges, relative information possible, still real hardware
            \item Time consuming (minutes to hours)
            \item Cheap (\$500-\$3k/board, ~\$300/license)
            \item Limited to process on board (~14-32nm)
        \end{itemize}
    \end{itemize}
    \centeredgraphicsmall{figures/fpga_vs_asic.png}
\end{frame}

\begin{frame}
    \frametitle{Need data, Part 2}
    \begin{itemize}
        \item Scoreboard great for integration testing, but not suitable for data collection
        \item Mini-scoreboard $+$ Deputy $+$ Wake $=$ Data
        \item xarrays \url{https://github.com/pydata/xarray}, N-D extension of pandas
        \item database vs netCDF \url{https://www.unidata.ucar.edu/blogs/developer/entry/netcdf_schema_language}
        \item Tensorflow(.js) to deploy trained models for front end design feedback
    \end{itemize}

    Typical gamma workload (lots of idle compute to be exploited):
    \centeredgraphic{figures/gamma_workload.png}
\end{frame}

\end{document}
